\documentclass[11pt]{article}
\usepackage[margin=1in]{geometry}
\usepackage{amsmath,amssymb,amsthm,mathtools,bm}
\usepackage{hyperref}
\hypersetup{colorlinks=true,linkcolor=blue,urlcolor=blue,citecolor=blue}

\title{Detailed Derivation Summary of \texttt{QAOAforSKp1/Basic.lean}}
\author{}
\date{\today}

\newcommand{\C}{\mathbb{C}}
\newcommand{\R}{\mathbb{R}}
\newcommand{\E}{\mathbb{E}}
\newcommand{\tr}{\operatorname{Tr}}
\newcommand{\ket}[1]{\lvert #1\rangle}
\newcommand{\bra}[1]{\langle #1\rvert}
\newcommand{\ip}[2]{\langle #1\mid #2\rangle}

\begin{document}
\maketitle

\section{Scope and Goal}
This note summarizes, in detail, the mathematical content currently formalized in
\texttt{QAOAforSKp1/Basic.lean}. The file develops:
\begin{enumerate}
  \item finite-dimensional state/operator infrastructure for \((n+1)\)-qubit systems,
  \item explicit computational-basis formulas for the mixer and SK cost operators,
  \item the QAOA generating function objective
  \[
    \E_J\!\left[\bra{s}U_C^\dagger U_B^\dagger e^{i\lambda C/(n+1)}U_B U_C\ket{s}\right],
  \]
  \item a sequence of algebraic rewrites (basis insertion, phase consolidation, variable changes, factorization),
  \item specialization to Gaussian couplings, giving explicit exponent structure and evaluation of the constant edge-sum term.
\end{enumerate}

\section{Hilbert-Space Model and Basic States}
\subsection{Indexing and states}
For each \(n\in\mathbb{N}\), the file uses
\[
  \texttt{BasisIdx}(n) = \mathrm{Fin}(2^{n+1}),
\]
so states are functions
\[
  \psi : \texttt{BasisIdx}(n) \to \C.
\]
The inner-product convention is encoded as
\[
  \texttt{bra}(\psi,\varphi) = \sum_{i\in\texttt{BasisIdx}(n)} \overline{\psi(i)}\,\varphi(i).
\]

\subsection{Single-qubit reference states}
The file defines computational basis states and \(\ket{+}\) from the Hadamard action on \(\ket{0}\), then proves
\[
  \ip{+}{+}=1.
\]

\subsection{Uniform product state \(\ket{s}\)}
The initial QAOA state is represented as the tensor-product uniform superposition over \((n+1)\) qubits:
\[
  \ket{s} = \ket{+}^{\otimes(n+1)}.
\]
Componentwise, each computational-basis amplitude is constant. The file proves
\[
  \ip{s}{s}=1.
\]
It also defines the pure density matrix
\[
  \rho_s = \ket{s}\!\bra{s}
\]
and proves
\[
  \tr(\rho_s)=1,
\]
including the identification \(\tr(\rho_s)=\ip{s}{s}\).

\section{Mixer Operator Infrastructure}
\subsection{Explicit kernel form}
The mixer is represented in computational basis by
\[
  (U_B\psi)(x) = \sum_y K_\beta(x,y)\,\psi(y),
\]
with kernel
\[
  K_\beta(x,y)=\cos(\beta)^{(n+1)-d(x,y)}\,\bigl(-i\sin\beta\bigr)^{d(x,y)},
\]
where \(d(x,y)\) is Hamming distance over \((n+1)\)-bit strings.

\subsection{Dagger and unitarity trace identity pathway}
The file includes the adjoint-kernel matrix expression and trace statements of the form
\[
  \tr\big((U_B^\dagger U_B)\rho_s\big)=1,
\]
under mixer-kernel orthonormality. A theorem placeholder remains:
\[
  \texttt{MixerKernelOrthonormal}\;\text{(currently with \texttt{sorry}).}
\]
This is the only remaining explicit proof gap warning in current builds.

\section{SK Cost Hamiltonian and Cost Unitary}
\subsection{Couplings and edge set}
Couplings are functions
\[
  J_{jk}:\mathrm{Fin}(n+1)\times\mathrm{Fin}(n+1)\to\R.
\]
The SK edge set is
\[
  \{(j,k): j<k\}.
\]
The file now proves exactly
\[
  |\{(j,k):j<k\}|=\frac{n(n+1)}{2}.
\]

\subsection{Spin map and SK cost on basis states}
For basis index \(z\), the spin value is
\[
  s_j(z)\in\{+1,-1\},
\]
via bit test (0\(\mapsto\)+1, 1\(\mapsto\)-1). The SK normalization is
\[
  \texttt{skNormFactor}(n)=\frac{1}{\sqrt{n+1}},
\]
with proved square identity
\[
  \texttt{skNormFactor}(n)^2=\frac{1}{n+1}.
\]
The basis cost is
\[
  C_J(z)=\frac{1}{\sqrt{n+1}}\sum_{j<k}J_{jk}s_j(z)s_k(z).
\]
Then \(U_C\) and \(U_C^\dagger\) are diagonal phases:
\[
  U_C(\gamma)\ket{z}=e^{-i\gamma C_J(z)}\ket{z},\qquad
  U_C(\gamma)^\dagger\ket{z}=e^{+i\gamma C_J(z)}\ket{z},
\]
with pointwise cancellation proved:
\[
  U_C^\dagger(z)U_C(z)=1.
\]

\section{Generating Function Objective}
The core object is
\[
  G_{n}(\beta,\gamma,\lambda)
  :=\E_J\!\left[\bra{s}U_C^\dagger U_B^\dagger e^{i\lambda C/(n+1)}U_B U_C\ket{s}\right].
\]
Here \(\lambda\in\R\) is a generating parameter (for later moment extraction via derivatives).

\section{Derivation Pipeline in the File}
This section tracks the theorem chain in the current Lean development.

\subsection{Step 1: Insert basis resolutions}
By expanding matrix trace and operator products in computational basis,
\[
  G_n = \E_J\!\left[\sum_{x,y,z}
    \bra{s}U_C^\dagger\ket{x}
    \bra{x}U_B^\dagger\ket{z}
    e^{i\lambda C_J(z)/(n+1)}
    \bra{z}U_B\ket{y}
    \bra{y}U_C\ket{s}
  \right].
\]
The file then substitutes explicit \(\ket{s}\)-entries and merges constants.

\subsection{Step 2: Consolidate cost phases}
Using diagonal form of \(U_C\), the product
\(
\bra{s}U_C^\dagger\ket{x}\,\bra{y}U_C\ket{s}
\)
becomes a single exponential with phase
\[
  i\gamma\,(C_J(x)-C_J(y)) + i\lambda\,\frac{C_J(z)}{n+1}.
\]

\subsection{Step 3: Swap integral and finite sums}
A theorem assumes integrability hypotheses needed to interchange expectation integral and finite sums; the expression is rewritten as finite sums of edge-integrals.

\subsection{Step 4: Pull out kernel factors}
Since mixer kernels do not depend on \(J\), they are taken outside the \(J\)-integral.

\subsection{Step 5: Spin-product variable change}
A change of variables corresponding to multiplying spin configurations by \(z\) is performed. The phase is rewritten in explicit edge form:
\[
  \exp\!\left(
    \frac{i}{\sqrt{n+1}}\sum_{j<k}J_{jk}\,\Theta_{jk}(x,y,z)
  \right),
\]
then reduced to a \(z=0\)-based coefficient representation under symmetry handling.

\subsection{Step 6: Edgewise symmetry and factorization assumptions}
Under explicit assumptions in theorems:
\begin{itemize}
  \item edgewise sign symmetry,
  \item edgewise factorization/independence of characteristic function contributions,
\end{itemize}
the expectation over \(J\) factors into edgewise characteristic functions
\[
  \chi_{jk}(t)=\E\big[e^{itJ_{jk}}\big].
\]
A notation
\[
  \phi_{jk}(x,y):=\gamma\,(x_jx_k-y_jy_k)
\]
is introduced, and the coefficient entering each edge-characteristic function is
\[
  a_{jk}(x,y)=\frac{1}{\sqrt{n+1}}\left(\phi_{jk}(x,y)+\frac{\lambda}{n+1}\right).
\]

\section{Gaussian Specialization}
Assuming each edge coupling is standard normal,
\[
  \chi_{jk}(t)=e^{-t^2/2}.
\]
This yields
\[
  \prod_{j<k} e^{-a_{jk}^2/2}
  = \exp\left(\sum_{j<k}-\frac{a_{jk}^2}{2}\right).
\]
The file then expands
\[
  \left(\phi_{jk}+\frac{\lambda}{n+1}\right)^2
  = \phi_{jk}^2 + 2\frac{\lambda}{n+1}\phi_{jk} + \left(\frac{\lambda}{n+1}\right)^2.
\]
So exponent is split into three edge-sums.

\section{Explicit Evaluation of the Constant Edge-Sum (Newest Step)}
The third sum is independent of \((j,k)\):
\[
  \sum_{j<k}-\frac12\cdot\frac{1}{n+1}\left(\frac{\lambda}{n+1}\right)^2
  = -\frac{\lambda^2}{2(n+1)^3}\,\binom{n+1}{2}
  = -\frac{n\lambda^2}{4(n+1)^2}.
\]
This is now derived in Lean (not assumed), using:
\begin{enumerate}
  \item exact cardinality proof for \(\texttt{skEdgeSet}\),
  \item exact identity \(\texttt{skNormFactor}(n)^2=1/(n+1)\).
\end{enumerate}

\section{Current Final Formula in the File}
The end-of-file comment and theorem chain give:
\begin{align*}
\E_J\!\left[\bra{s}U_C^\dagger U_B^\dagger e^{i\lambda C/(n+1)}U_B U_C\ket{s}\right]
&= \sum_{x,y}K_\beta^{(w)}(y)\,\overline{K_\beta^{(w)}(x)} \notag\\
\quad\times \exp\Bigg( &
-\frac{1}{2(n+1)} \sum_{j<k}\! \phi_{jk}(x,y)^2  \notag
-\frac{\lambda}{(n+1)^2}\sum_{j<k}\! \phi_{jk}(x,y)
-\frac{n\lambda^2}{4(n+1)^2}
\Bigg),
\end{align*}
with
\[
\phi_{jk}(x,y)=\gamma(x_jx_k-y_jy_k).
\]

\section{Assumptions vs. Derived Pieces}
\subsection*{Derived in code}
\begin{itemize}
  \item basis expansions and phase rewrites,
  \item Gaussian \(\prod\exp\to\exp\sum\) and square expansion,
  \item edge count \(|\texttt{skEdgeSet}|=n(n+1)/2\),
  \item explicit third-term evaluation \(-n\lambda^2/[4(n+1)^2]\),
  \item \(\texttt{skNormFactor}^2=1/(n+1)\).
\end{itemize}

\subsection*{Still assumed in theorems}
\begin{itemize}
  \item integrability assumption needed for exchanging expectation with finite sums,
  \item edgewise factorization/symmetry hypotheses used before Gaussian closure,
  \item mixer-kernel orthonormality theorem currently has a \texttt{sorry} placeholder.
\end{itemize}

\section{Practical Reading Map}
To inspect the Lean progression in order, follow theorem names around:
\begin{itemize}
  \item \texttt{ExpectedGeneratingFunction\_insert\_basis\_xyz}
  \item \texttt{ExpectedGeneratingFunction\_replace\_s\_factors}
  \item \texttt{ExpectedGeneratingFunction\_merge}
  \item \texttt{ExpectedGeneratingFunction\_swap\_integral\_and\_sum}
  \item \texttt{ExpectedGeneratingFunction\_pull\_out\_kernel\_factors}
  \item \texttt{ExpectedGeneratingFunction\_changeVariables\_xy\_spinProd}
  \item \texttt{ExpectedGeneratingFunction\_explicit\_phase\_after\_spinProd}
  \item \texttt{ExpectedGeneratingFunction\_edge\_indep\_and\_symm\_...}
  \item \texttt{ExpectedGeneratingFunction\_gaussian\_edges\_...}
  \item \texttt{ExpectedGeneratingFunction\_gaussian\_edges\_expSum\_thirdTerm\_evaluated}.
\end{itemize}

\end{document}
